\newcommand{\Rls}[1]{\ls{#1}_\mathbb{R}}

In \Cref{sec:comp:tor}, we define the $\mathbb{Q}$-divisor $D^T_L$ as a barycenter of the polytope. To find the barycenter of a moment polytope, we need to find a suitable measure on it.

As the character lattice $\ch(B) \cong \mathbb{Z}^{\rank(G)}$ is finite-dimensional, we can denote the closure of the moment ($\mathbb{Q}$-)polytope $P(X, L, \alpha) \subseteq \ch(B)_\mathbb{R} := \ch(B) \otimes_\mathbb{Z} \mathbb{R}$ by $\cl{P_\alpha}$. For every $m \in \mathbb{N}(L, \alpha)$, there is a measure $\dd \mu^\alpha_m$ on $\cl{P_\alpha}$, by
\begin{equation}
	\dd \mu^\alpha_m := {1 \over m^n} \sum_{\lambda \in P_m(X, L, \alpha)} N_{m, \lambda} \delta_{\lambda / m},
\end{equation}
where $\delta_u$ is the Dirac delta measure supported at the point $u \in \cl{P_\alpha}$, $n := \dim(X)$, and $N_{m, \lambda} = \dim(H^0(X, m L)_\lambda)$; $\dd \mu^\alpha_m$ is an atomic measure supported on ${1 \over m} P_m$.

\begin{lemdefn}
	There exists a measure $\dd \mu^\alpha$ on $\cl{P_\alpha}$, which is the weak limit of $\{\dd \mu^\alpha_m\}_{m = 1}^\infty$, and $\int_{\cl{P_\alpha}} \dd \mu^\alpha = \vol(L) / n!$.
\end{lemdefn}
\begin{remark*}
	It is the $G$-equivariant \emph{Duistermaat-Heckman measure} in \cite[Definition-Theorem 2.31]{Xu_Book}'s terminology. See also \cite{Delcroix_Kst_Fano_spher} and \cite[\S4.3]{Yin_valu_lin_sph}.
\end{remark*}
\begin{proof}
	See \cite[\S4.1]{Brion_grp_Pic_num_char_sph}.
\end{proof}

\begin{example}
	For toric varieties, the measure $\dd \mu$ is the lattice measure (or Lebesgue measure, uniform measure).
\end{example}

For a spherical $G$-variety $X$, there are finitely many $B$-invariant prime divisors, so the space of $B$-invariant $\mathbb{Q}$-Cartier $\mathbb{Q}$-divisors, denoted by $\CDiv(X)^B_\mathbb{Q}$, is a finite-dimensional $\mathbb{Q}$-vector space. Then, for every ample $\mathbb{Q}$-line bundle $L$, we can denote the closure of $B$-invariant $\mathbb{Q}$-linear series $\Qls{L}^B$ by $\Rls{L}^B \subseteq \CDiv(X)^B_\mathbb{R} := \CDiv(X)^B_\mathbb{Q} \otimes_\mathbb{Q} \mathbb{R}$.

Now, the extension $p_\alpha: \Rls{L}^B \to \cl{P_\alpha}$ is (still) a convex-linear isomorphism, and we can get the measure $\dd \mu_m$ and $\dd \mu$ on $\Rls{L}^B$.

\begin{proposition}\label[proposition]{prop:measONrls}
	The measures $\dd \mu_m$ (if they are well-defined) and $\dd \mu$ are independent of the choice of $\alpha$.
\end{proposition}
\begin{proof}
	We only need to prove $\dd \mu_m$ is independent of $\alpha$. Consider two $G$-linearization $\alpha, \beta$, and $m \in \mathbb{N}(L, \alpha) \cap \mathbb{N}(L, \beta)$, we have
	\begin{equation*}
		\phi := m \cdot (p_\beta \circ p_\alpha^{-1}) : P_m(X, L, \alpha) \to P_m(X, L, \beta).
	\end{equation*}
	For simplicity, define $G$-modules $V := (H^0(X, m L), \cdot_\alpha), W := (H^0(X, m L), \cdot_\beta)$, and the claim is equivalent to that for every $\lambda \in P_m(X, L, \alpha)$, the equality $\dim(V_\lambda) = \dim(W_{\phi(\lambda)})$ holds.
	
	To begin with, we have $V_\lambda^{(B)} = W_{\phi(\lambda)}^{(B)} = \mathbb{C} s$ for some $s \in H^0(X, m L)$. Let $D := \divQ(s)$, we can see $D \in {1 \over m} \ls{m L}^B$. For every $g \in G$,
	\begin{equation*}
		\divQ(g \cdot_\alpha s) = g \cdot D = \divQ(g \cdot_\beta s);
	\end{equation*}
	therefore, we get $V_\lambda = W_{\phi(\lambda)}$, a stronger result.
\end{proof}

% \begin{question}
% 	Is it possible to compute $\dd \mu$ without  $\alpha$?
% \end{question}

As the space $\Rls{L}^B$ is compact, the barycenter of measures $\dd \mu_m$ and $\dd \mu$ exists, denoted by $\bc_{\dd \mu_m}(\Rls{L}^B)$ and $\bc_{\dd \mu}(\Rls{L}^B)$, respectively. According to \cite[\S4.1]{Brion_grp_Pic_num_char_sph}, we can show $\bc_{\dd \mu}(\Rls{L}^B)$ is indeed a $\mathbb{Q}$-divisor.
